\section{Proposed Approach}

The proposed analysis synthesizes the retrieved circuit and computational literature around three circuit steps: MS GABAergic output, MEC inhibitory timing, CA1 temporal ensemble structure, and retrieval behavior. The first level treats medial septum GABAergic output as the upstream timing source. Work on theta-promoting septal circuits, optogenetic control, and sleep-wake state transitions establishes that septal activity can regulate hippocampal-entorhinal rhythmic state with sufficient temporal specificity to affect information routing \cite{arxiv151001100,arxiv221113591,arxiv250904454}. Because septal manipulation can also influence locomotion and arousal, those effects must be separated from retrieval-specific mechanisms \cite{arxiv181200625,arxiv210812261}.

The second level treats medial entorhinal inhibitory microcircuits as the transformation layer. The key question is whether septal GABAergic input changes the timing of local inhibitory populations that regulate entorhinal output. Parvalbumin-related and recurrent inhibitory motifs are relevant because they can narrow spike windows, normalize response distributions, and accelerate adaptation to changing input statistics \cite{arxiv160903622,arxiv220910634,arxiv240517745}. This level isolates timing as the controlled variable. Evidence that only firing rate or nonspecific excitability changes would not be sufficient.

The third level treats CA1 ensemble dynamics as the retrieval readout. CA1 should be evaluated through temporal structure, including phase relationships, cross-correlations, replay-like sequence organization, and ensemble composition. Retrieved models and analyses support the use of temporal coding, predictive coding, and calcium-imaging-based ensemble decoding as relevant readouts \cite{arxiv191211126,arxiv230511982,arxiv250809576}. The key test is whether altered entorhinal timing can change which CA1 temporal ensemble structure is expressed when prior and updated context-value associations can both be expressed under different retrieval conditions.

Support for the mechanism would require convergence across these levels. Septal control should alter medial entorhinal timing, medial entorhinal inhibitory control should reshape output precision, and CA1 should show retrieval-specific temporal reorganization. A dissociation at any level would constrain the model and redirect attention toward downstream output circuits, direct CA1 inhibition, or behavioral state variables.
